\documentclass[11pt,a4paper]{article}
\usepackage[utf8]{inputenc}
\usepackage[T1]{fontenc}
\usepackage[margin=24mm]{geometry}
\usepackage{mathpazo,booktabs,hyperref,microtype}
\hypersetup{colorlinks=true,urlcolor=blue}
\setlength{\parindent}{0pt}
\setlength{\parskip}{6pt}
\emergencystretch=2em
\begin{document}
\begin{center}
{\Large\bfseries Resumen técnico de decisiones y resultados}\\[4pt]
TITAN V4.5 --- revisión del manuscrito Life 4620846
\end{center}

Este resumen documenta las decisiones de la revisión a partir de sus archivos de evidencia, resultados y comprobaciones. Su alcance es la integración editorial y la verificación de agregados; no se ejecutaron nuevos entrenamientos ni inferencias de ECG durante esta revisión.

\section*{Decisiones y fundamento verificable}
\begin{enumerate}
\item \textbf{Conservar los resultados originales con sus denominadores.}
P6 mantiene 95,19\% de exactitud top-1 y 80,09\% de F1 binario medio en 1.643 ECG. P4 mantiene 80,23\% de exactitud binaria media y 79,35\% de F1 macro en 962 observaciones ECG--etiqueta de 733 ECG distintos. Son estimaciones retrospectivas seleccionadas; no se sustituyeron por las cifras del nuevo experimento.

\item \textbf{Autenticar el paquete recibido antes de integrarlo.}
El ZIP coincide byte por byte con el paquete público. El verificador comprobó 57 archivos y 46 registros de procedencia contra los objetos Git originales. La evidencia primaria está fijada al commit \texttt{63c74b3fde2c51f7f8cd}\allowbreak\texttt{e4ad396d77a1bd16cb29}. Los hashes permiten comprobar identidad de archivos; no certifican toda la historia de entrenamiento.

\item \textbf{Definir el alcance de CODE con precisión.}
El candidato recibió entrenamiento supervisado con pacientes de CODE y se evaluó en otros 3.891 pacientes reservados, con 4.921 ECG. La auditoría informa cero solapamiento de claves de paciente entre las cuatro particiones CODE. Esto permite describir una evaluación con pacientes separados dentro de CODE tras adaptación; la confirmación institucional independiente y la evaluación externa de P4 continúan pendientes. Los cinco objetivos son AFIB, SB, STACH, RBBB y 1AVB.

\item \textbf{Comparar los mismos ECG y conservar los resultados adversos.}
P6, P8 y el candidato usan los mismos ECG y referencias en esta comparación de cinco objetivos. Sus rutas de procesamiento y puntos de operación difieren, por lo que la ganancia no se atribuye a un componente aislado. Las pérdidas de sensibilidad AFIB/SB y el menor valor predictivo positivo de 1AVB frente a P8 se mantienen visibles. El candidato no superó el objetivo de F1 macro 0,70 ni todas las salvaguardas de promoción.
\end{enumerate}

\section*{Resultados del experimento CODE}
\begin{center}
\small
\begin{tabular}{lrrr}
\toprule
Sistema & F1 macro [IC 95\%] & AUROC macro & AP macro\\
\midrule
P6 congelado & 34,80 [32,00--37,50] & 92,40 & 35,74\\
P8 congelado & 37,04 [33,96--40,04] & 93,11 & 39,00\\
Candidato adaptado & 60,27 [56,46--63,41] & 97,81 & 57,95\\
\bottomrule
\end{tabular}
\end{center}
Las cifras de la tabla son porcentajes. Los intervalos de F1 usan 1.000 remuestreos por paciente; no se proporcionaron intervalos macro de AUROC/AP. Las diferencias pareadas de F1 son +25,47 puntos porcentuales frente a P6 (IC 95\%: +22,20 a +28,83) y +23,23 frente a P8 (+19,77 a +26,87). El candidato omitió 45 de 101 ECG positivos para AFIB y 30 de 64 positivos para SB. Su exactitud binaria media de 98,31\% no equivale a exactitud diagnóstica completa por ECG.

\section*{Qué se comprobó y dónde contrastarlo}
Se recalcularon métricas derivadas de matrices de confusión y medias aritméticas públicas. AUROC, AP e intervalos bootstrap proceden de agregados autenticados mediante recibos; no se recalcularon localmente desde predicciones privadas. El artículo quedó en 45 páginas, con cuatro tablas y dos figuras adicionales. La extracción limpia del ZIP compiló al mismo texto PDF normalizado; pasaron 33 comprobaciones locales, Ruff y cinco pruebas dirigidas.

Archivos de trazabilidad, con rutas relativas al directorio de la revisión:
\begin{itemize}
\item \path{audit/CODE_EVIDENCE_INTEGRATION.json}
\item \path{supplementary/MANUSCRIPT_EVIDENCE_LEDGER.csv}
\item \path{FINAL_MANUSCRIPT_SCIENTIFIC_QA.json}
\end{itemize}
Las tablas CODE completas están en los CSV \path{supplementary/CODE_*}; el procedimiento y las limitaciones se desarrollan en \path{latex/code_methods.tex}, \path{latex/code_results.tex} y \path{latex/code_discussion.tex}. La documentación conserva las limitaciones sobre procedencia histórica y sobre los informes originales incompletos de los revisores.

\href{https://github.com/Mateo197802/titan-v45-ecg-multihead/pull/23}{PR borrador \#23: fuentes, resultados y documentación de la revisión.}
\end{document}
